% Chapter 4 -- Model Design and Training
% Figures are authored in TikZ so the chapter compiles without external assets.
% Placeholders [CITATION: ...] mark references still to be added.

\label{ch:method}

\tikzset{
  blk/.style={draw, rounded corners=2pt, align=center, inner sep=4pt,
              minimum height=8mm, font=\footnotesize},
  grp/.style={draw, dashed, rounded corners=3pt, inner sep=6pt},
  flow/.style={-{Latex[length=2mm]}, thick},
}

This chapter describes the models used for activity recognition and the
procedure used to train and evaluate them. It begins with an overview of the
experimental design and its relation to the research questions. It then covers
the input representations derived from the inertial and acoustic windows of
Chapter~\ref{ch:data}, the single-modality models, the three fusion strategies,
and a classical baseline using hand-crafted features. It closes with the training
procedure and the evaluation protocol that all models share. The results obtained
with these models are presented in Chapter~\ref{ch:results}, and the conversion
of selected models for execution on the wearable device is covered in
Chapter~\ref{ch:implementation}.

\section{Overview of the Experimental Design}
\label{sec:design}

The experiments compare models that differ in one respect at a time while
sharing everything else: the same windows, the same train/validation split, the
same input normalisation and the same training settings. This makes differences
in accuracy attributable to the model design rather than to differences in the
data or in the training procedure.

The models fall into four groups, each linked to one or more of the research
questions formulated in Chapter~1:

\begin{itemize}
  \item \textbf{Single-modality models} use either inertial or acoustic data
  alone. They form the reference point against which the benefit of adding audio
  is measured (RQ1, RQ3).
  \item \textbf{Fusion models} combine both modalities at different stages of
  processing: at the input level (early fusion), at the level of learned
  representations (embedding-level fusion), or at the level of class predictions
  (decision-level fusion). Their comparison addresses RQ2.
  \item \textbf{A classical baseline} based on statistical features and a random
  forest shows how much the learned models gain over a conventional approach
  (RQ3).
  \item \textbf{Reduced model variants} with fewer parameters were designed with
  the memory and computational limits of the target device in mind. They connect
  the offline comparison to the deployment study in
  Chapter~\ref{ch:implementation} (RQ4).
\end{itemize}

\begin{table}[htbp]
\centering
\small
\renewcommand{\arraystretch}{1.2}
\begin{tabularx}{\textwidth}{|p{3.1cm}|X|p{2.4cm}|p{2.6cm}|}
\hline
\textbf{Group} & \textbf{Model} & \textbf{Inputs} & \textbf{Fusion stage} \\
\hline
Single modality & IMU CNN & inertial & -- \\
Single modality & Audio CNN (GAP, GRU and hybrid pooling) & acoustic & -- \\
Early fusion & Channel concatenation (standard and small) & inertial + acoustic & input \\
Early fusion & Unified spectrogram & inertial + acoustic & input \\
Embedding-level fusion & Frozen backbones with a trained fusion head (standard and small) & inertial + acoustic & representation \\
Decision-level fusion & Weighted average with global or per-class weights (standard and small) & inertial + acoustic & prediction \\
Classical baseline & Random forest on hand-crafted features & inertial + acoustic & feature \\
\hline
\end{tabularx}
\caption{Overview of the evaluated models.}
\label{tab:model-overview}
\end{table}

All experiments use the dominant-wrist dataset described in
Chapter~\ref{ch:data}, with 2565 paired windows across nine activity classes. The
standalone audio model is the one exception: it additionally uses audio windows
from the non-dominant wrist and the waist, as explained in
Section~\ref{sec:audio-model}.

\section{Input Representations}
\label{sec:inputs}

\subsection{Inertial Windows}
\label{sec:imu-input}

Each inertial window holds 160 samples of the six channels (three-axis
acceleration and three-axis angular velocity) recorded at 80\,Hz. Before entering
a model, the window is resampled to 200 time steps by linear interpolation along
the time axis, so that all inertial models take inputs of size $200 \times 6$.
The target length of 200 was chosen to match the time resolution of the audio
spectrogram, which has 198 frames per window (Section~\ref{sec:logmel}). This
allows both modalities to share one time axis in the early fusion model
(Section~\ref{sec:early-fusion}) without coarsening the audio representation, and
using the same length in all inertial models keeps their inputs identical across
the experiments. Missing or non-finite values, which can remain at the very edges of a recording
after decimation, are replaced by zero.

Accelerometer and gyroscope channels have different physical units and value
ranges. Each channel is therefore standardised to zero mean and unit variance.
The mean and standard deviation of each channel are computed on the training
split only and then applied unchanged to the validation data, so that no
statistics of the validation set enter the training process. The same statistics
are later embedded in the device firmware, where they are applied to the live
sensor data (Chapter~\ref{ch:implementation}).

\subsection{Audio Log-Mel Spectrograms}
\label{sec:logmel}

Raw audio windows of 32\,000 samples are too long to be used directly as input
to a small model. Each window is therefore converted into a log-mel spectrogram,
a time-frequency representation that is widely used for sound classification
{[}CITATION{]}. The conversion consists of three steps:

\begin{enumerate}
  \item a short-time Fourier transform (STFT) with frames of 25\,ms (400
  samples), a hop of 10\,ms (160 samples) and an FFT length of 512;
  \item a projection of the magnitude spectrum onto 40 mel-spaced frequency bands
  between 20\,Hz and 8\,kHz {[}CITATION{]};
  \item the natural logarithm of the band energies, with a small constant
  ($10^{-6}$) added to avoid the logarithm of zero.
\end{enumerate}

A two-second window produces a spectrogram of 198 frames $\times$ 40 mel bands.
This reduces the input size by roughly a factor of four compared with the raw
waveform, while keeping the spectral and temporal structure that distinguishes,
for example, the continuous whine of a powered screwdriver from the short,
broadband impacts of hammering. As with the inertial data, the spectrogram is
standardised per mel band using the mean and standard deviation of the training
split. Figure~\ref{fig:logmel} summarises the feature extraction, and
Table~\ref{tab:inputs} lists the resulting input representations.

\begin{figure}[htbp]
\centering
\begin{tikzpicture}[node distance=3.5mm, every node/.style={blk, text width=7.2cm}]
  \node (wav) {Audio window: 2\,s, 32\,000 samples at 16\,kHz};
  \node (stft) [below=of wav] {STFT: 25\,ms frames, 10\,ms hop, FFT length 512};
  \node (mag) [below=of stft] {Magnitude spectrum (257 frequency bins per frame)};
  \node (mel) [below=of mag] {Mel filter bank: 40 bands, 20\,Hz--8\,kHz};
  \node (log) [below=of mel] {Natural logarithm, $\log(E + 10^{-6})$};
  \node (norm) [below=of log] {Per-band standardisation (training statistics)};
  \node (out) [below=of norm] {Log-mel spectrogram: 198 frames $\times$ 40 bands};
  \foreach \a/\b in {wav/stft, stft/mag, mag/mel, mel/log, log/norm, norm/out}
    \draw[flow] (\a) -- (\b);
\end{tikzpicture}
\caption{Conversion of a two-second audio window into the log-mel spectrogram used
as input to the acoustic models.}
\label{fig:logmel}
\end{figure}

\begin{table}[htbp]
\centering
\small
\renewcommand{\arraystretch}{1.2}
\begin{tabularx}{\textwidth}{|l|p{3.1cm}|p{3.3cm}|X|}
\hline
\textbf{Modality} & \textbf{Raw window} & \textbf{Model input} & \textbf{Normalisation} \\
\hline
Inertial & $160 \times 6$ at 80\,Hz & $200 \times 6$ (linear resampling) & per channel, training statistics \\
Audio & 32\,000 samples at 16\,kHz & $198 \times 40$ log-mel & per mel band, training statistics \\
\hline
\end{tabularx}
\caption{Input representations of the two modalities.}
\label{tab:inputs}
\end{table}

\section{Single-Modality Models}
\label{sec:single}

\subsection{Inertial Model}
\label{sec:imu-model}

The inertial model is a one-dimensional convolutional neural network (1D CNN)
that convolves along the time axis of the window. 1D CNNs are a common choice for
wearable activity recognition because they learn local motion patterns directly
from the raw signal and remain small enough for embedded use {[}CITATION{]}.

The network has two convolutional blocks. The first block applies 64 filters with
a kernel size of 5, followed by batch normalisation and max pooling with a factor
of 2. The second block applies 128 filters with a kernel size of 3, followed by
batch normalisation. Global average pooling then reduces the temporal dimension
to a single feature vector. This vector passes through a dense layer with 64
units, which also serves as the embedding for embedding-level fusion
(Section~\ref{sec:embedding-fusion}), and through dropout with a rate of 0.3,
before a softmax output layer with one unit per class.

\subsection{Audio Model and Temporal Pooling Variants}
\label{sec:audio-model}

The audio model treats the log-mel spectrogram as a single-channel image and
processes it with a two-dimensional CNN {[}CITATION{]}. It has two convolutional
blocks with 32 and 64 filters of size $3 \times 3$, each followed by batch
normalisation and $2 \times 2$ max pooling. The resulting feature map has a size
of $49 \times 10 \times 64$ (time $\times$ frequency $\times$ channels).

How this feature map is reduced to a single vector determines which properties of
the sound the classifier can use. Three variants were therefore evaluated
(Figure~\ref{fig:audio-variants}):

\begin{itemize}
  \item \textbf{Global average pooling (GAP).} The feature map is averaged over
  both time and frequency {[}CITATION{]}. This is the standard approach and gives a compact model. However,
  averaging over time removes the information about how the sound develops within
  the window. Two activities with a similar average spectrum but a different
  temporal pattern, such as the regular strikes of a hammer and the steady sound
  of a screwdriver, can then become hard to separate.
  \item \textbf{Recurrent pooling (GRU).} The feature map is reshaped into a
  sequence of 49 time steps with 640 features each and passed through a
  bidirectional gated recurrent unit (GRU) with 32 units per direction
  {[}CITATION{]}. The GRU summarises the sequence while keeping its
  temporal order, so the model can use rhythm and envelope information.
  \item \textbf{Hybrid pooling.} Both summaries are computed in parallel and
  concatenated: the GAP vector, which describes the average timbre of the
  window, and the GRU output, which describes its temporal structure. The
  classifier can then rely on whichever information is useful for a given class.
\end{itemize}

In all three variants, the pooled vector passes through a dense layer with 64
units and dropout with a rate of 0.4, followed by the softmax output.

\begin{figure}[htbp]
\centering
\resizebox{\textwidth}{!}{%
\begin{tikzpicture}[node distance=3.5mm and 6mm, every node/.style={blk}]
  \node[text width=3.4cm] (in) {Input\\$198 \times 40 \times 1$};
  \node[text width=3.4cm, below=of in] (c1) {Conv2D 32, $3\times3$\\BN, MaxPool $2\times2$};
  \node[text width=3.4cm, below=of c1] (c2) {Conv2D 64, $3\times3$\\BN, MaxPool $2\times2$};
  \node[text width=3.4cm, below=of c2] (fm) {Feature map\\$49 \times 10 \times 64$};
  \foreach \a/\b in {in/c1, c1/c2, c2/fm} \draw[flow] (\a) -- (\b);

  % (a) GAP
  \node[text width=2.8cm, right=12mm of in] (ga) {(a) GAP};
  \node[text width=2.8cm, below=of ga] (gd) {Dense 64\\Dropout 0.4};
  \node[text width=2.8cm, below=of gd] (go) {Softmax (9)};
  \draw[flow] (ga) -- (gd); \draw[flow] (gd) -- (go);

  % (b) GRU
  \node[text width=2.8cm, right=of ga] (ra) {(b) Reshape $49 \times 640$};
  \node[text width=2.8cm, below=of ra] (rg) {BiGRU, 32 units};
  \node[text width=2.8cm, below=of rg] (rd) {Dense 64\\Dropout 0.4};
  \node[text width=2.8cm, below=of rd] (ro) {Softmax (9)};
  \foreach \a/\b in {ra/rg, rg/rd, rd/ro} \draw[flow] (\a) -- (\b);

  % (c) hybrid
  \node[text width=2.8cm, right=of ra] (ha) {(c) GAP $\parallel$ BiGRU};
  \node[text width=2.8cm, below=of ha] (hc) {Concatenate\\(64 + 64)};
  \node[text width=2.8cm, below=of hc] (hd) {Dense 64\\Dropout 0.4};
  \node[text width=2.8cm, below=of hd] (ho) {Softmax (9)};
  \foreach \a/\b in {ha/hc, hc/hd, hd/ho} \draw[flow] (\a) -- (\b);

  \draw[flow] (fm.east) -| ($(ga.west)+(-5mm,0)$) -- (ga.west);
  \draw[flow] (fm.east) -- ++(4mm,0) |- ($(ra.north)+(0,4mm)$) -- (ra.north);
  \draw[flow] (fm.east) -- ++(4mm,0) |- ($(ha.north)+(0,4mm)$) -- (ha.north);
\end{tikzpicture}}
\caption{Audio model with a shared convolutional stem and the three pooling
variants: (a) global average pooling, (b) recurrent pooling with a bidirectional
GRU, and (c) hybrid pooling, which concatenates both summaries.}
\label{fig:audio-variants}
\end{figure}

Because audio content depends less on where the sensor is worn than inertial
data does, the standalone audio model is trained on audio windows from three body
positions: the dominant wrist, the non-dominant wrist and the waist. This roughly
triples the amount of acoustic training data. The pooled audio is used only for
the standalone audio model. All fusion models use dominant-wrist audio
exclusively, so that every audio window is paired with the inertial window from
the same device and time interval.

The pooling variants are compared on the standalone audio model. Within the
fusion models, the acoustic component uses the GAP variant. With 23\,945
parameters it is about six times smaller than the recurrent variants, and it
avoids recurrent layers, whose sequential computation is costly on the target
microcontroller. This keeps all fusion models within reach of the deployment study
in Chapter~\ref{ch:implementation}. The performance of the three variants is
reported in Chapter~\ref{ch:results}. Table~\ref{tab:single-arch} summarises both
single-modality architectures.

\begin{table}[htbp]
\centering
\small
\renewcommand{\arraystretch}{1.2}
\begin{tabularx}{\textwidth}{|l|X|X|}
\hline
\textbf{Layer} & \textbf{IMU CNN} & \textbf{Audio CNN} \\
\hline
Input & $200 \times 6$ & $198 \times 40 \times 1$ \\
Block 1 & Conv1D 64, kernel 5, BN, MaxPool 2 & Conv2D 32, $3 \times 3$, BN, MaxPool $2 \times 2$ \\
Block 2 & Conv1D 128, kernel 3, BN & Conv2D 64, $3 \times 3$, BN, MaxPool $2 \times 2$ \\
Pooling & Global average & GAP / BiGRU(32) / GAP + BiGRU \\
Head & Dense 64, Dropout 0.3 & Dense 64, Dropout 0.4 \\
Output & Softmax (9 classes) & Softmax (9 classes) \\
\hline
Parameters & 36\,297 & 23\,945 (GAP) / 153\,353 (GRU) / 157\,449 (hybrid) \\
\hline
\end{tabularx}
\caption{Architecture of the single-modality models. All convolutional and dense
layers except the output use ReLU activations.}
\label{tab:single-arch}
\end{table}

\section{Fusion Architectures}
\label{sec:fusion}

The three fusion strategies differ in the point along the processing chain at
which the two modalities are combined {[}CITATION{]}.
Figure~\ref{fig:fusion-overview} gives an overview.

\begin{figure}[htbp]
\centering
\resizebox{\textwidth}{!}{%
\begin{tikzpicture}[node distance=4mm, every node/.style={blk}]
  % --- early fusion ---
  \node[text width=1.9cm] (e-imu) {IMU\\$200 \times 6$};
  \node[text width=1.9cm, right=3mm of e-imu] (e-aud) {Log-mel\\$\to 200 \times 40$};
  \node[text width=4.3cm, below=8mm of $(e-imu)!0.5!(e-aud)$] (e-cat) {Concatenate\\$200 \times 46$};
  \node[text width=4.3cm, below=of e-cat] (e-cnn) {1D CNN};
  \node[text width=4.3cm, below=of e-cnn] (e-out) {Softmax};
  \draw[flow] (e-imu) -- (e-cat); \draw[flow] (e-aud) -- (e-cat);
  \draw[flow] (e-cat) -- (e-cnn); \draw[flow] (e-cnn) -- (e-out);
  \node[draw=none, above=5mm of $(e-imu)!0.5!(e-aud)$, font=\small\bfseries] {Early fusion};

  % --- embedding-level fusion ---
  \node[text width=1.9cm, right=14mm of e-aud] (m-imu) {IMU};
  \node[text width=1.9cm, right=3mm of m-imu] (m-aud) {Log-mel};
  \node[text width=1.9cm, below=of m-imu] (m-icnn) {IMU CNN\\(frozen)};
  \node[text width=1.9cm, below=of m-aud] (m-acnn) {Audio CNN\\(frozen)};
  \node[text width=4.3cm, below=8mm of $(m-icnn)!0.5!(m-acnn)$] (m-cat) {Concatenate embeddings\\$64 + 64 = 128$};
  \node[text width=4.3cm, below=of m-cat] (m-head) {Fusion head\\Dense 64, Dropout 0.4};
  \node[text width=4.3cm, below=of m-head] (m-out) {Softmax};
  \draw[flow] (m-imu) -- (m-icnn); \draw[flow] (m-aud) -- (m-acnn);
  \draw[flow] (m-icnn) -- (m-cat); \draw[flow] (m-acnn) -- (m-cat);
  \draw[flow] (m-cat) -- (m-head); \draw[flow] (m-head) -- (m-out);
  \node[draw=none, above=5mm of $(m-imu)!0.5!(m-aud)$, font=\small\bfseries] {Embedding-level fusion};

  % --- decision-level fusion ---
  \node[text width=1.9cm, right=14mm of m-aud] (d-imu) {IMU};
  \node[text width=1.9cm, right=3mm of d-imu] (d-aud) {Log-mel};
  \node[text width=1.9cm, below=of d-imu] (d-icnn) {IMU CNN};
  \node[text width=1.9cm, below=of d-aud] (d-acnn) {Audio CNN};
  \node[text width=1.9cm, below=of d-icnn] (d-ip) {$\mathbf{p}^{\mathrm{IMU}}$};
  \node[text width=1.9cm, below=of d-acnn] (d-ap) {$\mathbf{p}^{\mathrm{audio}}$};
  \node[text width=4.3cm, below=8mm of $(d-ip)!0.5!(d-ap)$] (d-avg) {Weighted average\\(global or per-class weights)};
  \node[text width=4.3cm, below=of d-avg] (d-out) {Argmax};
  \draw[flow] (d-imu) -- (d-icnn); \draw[flow] (d-aud) -- (d-acnn);
  \draw[flow] (d-icnn) -- (d-ip); \draw[flow] (d-acnn) -- (d-ap);
  \draw[flow] (d-ip) -- (d-avg); \draw[flow] (d-ap) -- (d-avg);
  \draw[flow] (d-avg) -- (d-out);
  \node[draw=none, above=5mm of $(d-imu)!0.5!(d-aud)$, font=\small\bfseries] {Decision-level fusion};
\end{tikzpicture}}
\caption{The three fusion strategies. Early fusion combines the raw inputs before
learning, embedding-level fusion combines learned 64-dimensional representations
of two frozen backbones, and decision-level fusion combines the class
probabilities of two independently trained models.}
\label{fig:fusion-overview}
\end{figure}

\subsection{Early Fusion}
\label{sec:early-fusion}

Early fusion combines the two modalities before any learning takes place, so
that a single network learns from both signals jointly. Two forms were
implemented.

\textbf{Channel concatenation.} The log-mel spectrogram ($198 \times 40$) is
resampled along its time axis to 200 frames, matching the length of the inertial
window. The two inputs are then concatenated along the feature axis, giving a
$200 \times 46$ input with 6 inertial and 40 acoustic channels. Each of the 46
channels is standardised with its own training statistics. This step matters more
here than in the single-modality models, because raw inertial values and log-mel
energies lie on very different scales. The fused input is processed by a 1D CNN
with the same structure as the inertial model (Section~\ref{sec:imu-model}).

\textbf{Unified spectrogram.} Here the inertial data are also converted into
spectrograms, so that both modalities share one image-like representation. A
short-time Fourier transform with a frame length of 20 samples, a hop of 2
samples and an FFT length of 32 is applied to each of the six inertial axes. The
log magnitudes are resized to the $198 \times 40$ grid of the audio spectrogram
by bilinear interpolation. The six inertial spectrograms and the audio
spectrogram are stacked as the seven channels of a single input and processed by
the 2D audio CNN with global average pooling.

\subsection{Embedding-Level Fusion}
\label{sec:embedding-fusion}

Embedding-level fusion first trains the inertial and the acoustic model
separately, as described in Section~\ref{sec:single}. Their weights are then
frozen, and for each window the outputs of the 64-unit dense layer are taken as
embeddings. The two 64-dimensional embeddings are concatenated into a
128-dimensional vector, which is the input to a small fusion head: a dense layer
with 64 units, dropout with a rate of 0.4 and a softmax output (8841 parameters).
Only the fusion head is trained in this second stage. It can learn how the two
modalities relate to each other without changing the feature extraction that each
backbone has already learned.

\subsection{Decision-Level Fusion}
\label{sec:decision-fusion}

Decision-level fusion combines the class probabilities predicted by the two
independently trained models {[}CITATION{]}. For a window $x$, the fused probability
of class $c$ is a weighted average
\begin{equation}
  p_c(x) = w \cdot p^{\mathrm{IMU}}_c(x) + (1 - w) \cdot p^{\mathrm{audio}}_c(x),
  \qquad w \in [0, 1],
  \label{eq:global-fusion}
\end{equation}
and the predicted class is the one with the highest fused probability. The
weight $w$ is selected by a grid search over the values $0, 0.05, \dots, 1$,
keeping the value that gives the highest accuracy on the validation data.
Compared with a trained fusion head, this method has only a single free
parameter and is therefore less prone to overfitting on a small dataset.

A single weight, however, assumes that the relative reliability of the two
modalities is the same for every class. This assumption does not hold for the
present activity set. For some activities the motion pattern is distinctive
while the sound is not; standing idle and measuring, for example, produce hardly
any sound. For others the opposite is true, and the powered screwdriver is the
clearest example of an activity with a distinctive acoustic signature. The method
was therefore extended to one weight per class, which gives a weight vector
$\mathbf{w} = (w_1, \dots, w_C)$ and
\begin{equation}
  p_c(x) = w_c \cdot p^{\mathrm{IMU}}_c(x) + (1 - w_c) \cdot p^{\mathrm{audio}}_c(x).
  \label{eq:perclass-fusion}
\end{equation}
The only difference to Equation~\ref{eq:global-fusion} is that the weight now
carries the index $c$. When the fused model computes the score of a particular
class, it uses the weight of that class, and it can therefore rely on whichever
sensor is more informative for that class. A value of $w_c$ close to 1 means that
the score of class $c$ is determined mainly by the inertial model, and a value
close to 0 means that it is determined mainly by the acoustic model.

A short example illustrates the effect. Consider a window of the class
\texttt{measuring}, for which the two models produce the probabilities shown in
Table~\ref{tab:fusion-example}; only three classes are shown for readability. The
inertial model recognises the activity, whereas the acoustic model hears almost
silence and leans towards \texttt{idle}. With a global weight of $w = 0.5$, the
correct class wins only by a small margin. If the weight of \texttt{measuring} is
raised to $w_{\mathrm{measuring}} = 0.9$, its score is taken almost entirely from
the more reliable inertial model, and the margin increases clearly. The score of
\texttt{idle} does not change, because its own weight has not changed.

\begin{table}[htbp]
\centering
\small
\renewcommand{\arraystretch}{1.2}
\begin{tabular}{|l|c|c|c|}
\hline
 & \texttt{measuring} & \texttt{idle} & \texttt{hammering} \\
\hline
$\mathbf{p}^{\mathrm{IMU}}$ & 0.70 & 0.25 & 0.05 \\
$\mathbf{p}^{\mathrm{audio}}$ & 0.20 & 0.60 & 0.20 \\
\hline
Global weight, $w = 0.5$ & \textbf{0.45} & 0.425 & 0.125 \\
Per-class weight, $w_{\mathrm{measuring}} = 0.9$, others 0.5 & \textbf{0.65} & 0.425 & 0.125 \\
\hline
\end{tabular}
\caption{Illustrative example of global and per-class decision-level fusion for a
\texttt{measuring} window. The values are chosen for illustration and are not
taken from the experiments.}
\label{tab:fusion-example}
\end{table}

The nine weights cannot be found by simply trying every combination. With 21
candidate values per weight, this would require $21^9 \approx 8 \cdot 10^{11}$
evaluations. Instead, the weights are determined by coordinate ascent, which
adjusts one weight at a time while keeping the others fixed:

\begin{enumerate}
  \item \textbf{Initialisation.} All nine class weights are set to the best global
  weight from Equation~\ref{eq:global-fusion}. The search therefore starts from
  the result of global fusion and cannot end with a lower validation accuracy.
  \item \textbf{Adjusting one class.} For the first class, its weight is set in
  turn to each of the 21 values $0, 0.05, \dots, 1$, while the weights of all
  other classes stay unchanged. For each value, all validation windows are fused
  with Equation~\ref{eq:perclass-fusion} and the overall accuracy is computed. The
  weight is replaced only if one of the values gives a strictly higher accuracy
  than the current setting.
  \item \textbf{All classes.} Step~2 is repeated for the second class, the third
  class and so on, until every class has been adjusted once. This is one pass.
  \item \textbf{Second pass.} Because changing one weight can slightly shift the
  best value of another, the whole pass is repeated once more.
\end{enumerate}

In total, the search evaluates $2 \times 9 \times 21 = 378$ weight settings, which
takes only a fraction of a second. Two passes are sufficient because the classes
interact only weakly. Changing $w_c$ affects only those windows whose decision was
already close between class $c$ and another class; for most windows one class
leads clearly, and a single weight cannot change the outcome. The resulting
weights can be read as the degree of trust that the fused model places in each
modality for each class, which also makes them useful for interpreting the
results in Chapter~\ref{ch:results}. Figure~\ref{fig:perclass-fusion} summarises
the procedure.

\begin{figure}[htbp]
\centering
\begin{tikzpicture}[node distance=4mm and 6mm, every node/.style={blk}]
  \node[text width=2.6cm] (pi) {$\mathbf{p}^{\mathrm{IMU}}$\\(9 probabilities)};
  \node[text width=2.6cm, right=of pi] (w) {$\mathbf{w}$\\(9 class weights)};
  \node[text width=2.6cm, right=of w] (pa) {$\mathbf{p}^{\mathrm{audio}}$\\(9 probabilities)};
  \node[text width=8.4cm, below=7mm of w] (comb) {$p_c = w_c\, p^{\mathrm{IMU}}_c + (1 - w_c)\, p^{\mathrm{audio}}_c$ for each class $c$};
  \node[text width=8.4cm, below=of comb] (am) {Predicted class $= \arg\max_c p_c$};
  \draw[flow] (pi) -- (comb); \draw[flow] (w) -- (comb); \draw[flow] (pa) -- (comb);
  \draw[flow] (comb) -- (am);

  \node[text width=8.4cm, below=10mm of am] (s1) {Initialise all $w_c$ with the best global weight};
  \node[text width=8.4cm, below=of s1] (s2) {For each class $c$: scan $w_c \in \{0, 0.05, \dots, 1\}$, keep the value with the highest validation accuracy};
  \node[text width=8.4cm, below=of s2] (s3) {Repeat for two passes over all classes};
  \draw[flow] (s1) -- (s2); \draw[flow] (s2) -- (s3);
  \node[grp, fit=(s1)(s2)(s3), label={[font=\footnotesize]above:Weight search (coordinate ascent)}] {};
\end{tikzpicture}
\caption{Decision-level fusion with per-class weights (top) and the coordinate-ascent
search used to determine the weights on the validation data (bottom).}
\label{fig:perclass-fusion}
\end{figure}

\section{Classical Baseline}
\label{sec:classical}

To place the learned models in context, a conventional approach based on
hand-crafted features was also evaluated. From each inertial window, six
statistics are computed per channel: mean, standard deviation, minimum, maximum,
root mean square and zero-crossing rate, which gives 36 inertial features. From
each log-mel spectrogram, the mean and standard deviation of every mel band over
time are computed, which gives 80 acoustic features. The two feature sets are
concatenated into a 116-dimensional vector (Table~\ref{tab:features}), so the
baseline also uses feature-level fusion.

A random forest with 300 trees is trained on these vectors, with balanced class
weights to account for the class imbalance {[}CITATION{]}. Random forests need no feature scaling and give robust results on small
tabular datasets, which makes them a common reference point in activity
recognition {[}CITATION{]}.

\begin{table}[htbp]
\centering
\small
\renewcommand{\arraystretch}{1.2}
\begin{tabularx}{\textwidth}{|l|X|r|}
\hline
\textbf{Modality} & \textbf{Features} & \textbf{Count} \\
\hline
Inertial & mean, standard deviation, minimum, maximum, RMS and zero-crossing rate for each of the 6 channels & 36 \\
Audio & mean and standard deviation over time for each of the 40 mel bands & 80 \\
\hline
\multicolumn{2}{|l|}{\textbf{Total}} & \textbf{116} \\
\hline
\end{tabularx}
\caption{Hand-crafted features of the classical baseline.}
\label{tab:features}
\end{table}

\section{Deployment-Oriented Model Variants}
\label{sec:small-models}

The models described so far were designed for offline evaluation. On the target
microcontroller, however, inference time and memory usage are strict limits.
The TensorFlow Lite Micro build for this platform executes convolutions with
unaccelerated reference kernels, so their cost grows directly with the number of
filters and with the size of the feature maps (Chapter~\ref{ch:implementation}).
Reduced
variants were therefore created that keep the overall architecture but use fewer
filters:

\begin{itemize}
  \item \textbf{Small early fusion model.} The structure of the
  channel-concatenation model with 32 and 64 instead of 64 and 128 filters. This
  reduces the number of multiply-accumulate operations, and with it the inference
  time, by roughly a factor of four.
  \item \textbf{Small inertial backbone.} A 1D CNN with 32 and 64 filters.
  \item \textbf{Small audio backbone.} A 2D CNN with 16 and 32 filters. Its
  convolutions use strides of 4 and 2 instead of separate pooling layers, so the
  downsampling happens within the convolution itself. As a result, a feature map
  at the full $198 \times 40$ resolution never has to be held in memory, which
  would otherwise take up a large part of the device's working memory.
\end{itemize}

The small backbones are used for the small embedding-level and decision-level
fusion models. Because these models have less capacity, three regularisation
measures were added to reduce overfitting: L2 weight decay ($10^{-4}$) on all
learned layers, higher dropout rates (0.4 and 0.5), and label smoothing of 0.1 in
the loss function {[}CITATION{]}.

In addition, the training data of the small backbones are augmented. One
augmented copy is generated for each training window, so that the training set
doubles in size; the validation data are never augmented.

\begin{itemize}
  \item \textbf{Inertial windows} receive a small random three-dimensional rotation
  of up to $\pm 15^\circ$, applied jointly to the accelerometer and gyroscope axes
  to mimic small changes in how the device sits on the wrist, together with
  per-channel scaling (factor 0.9--1.1), additive noise (3\,\% of the channel's
  standard deviation) and a mild time warp (speed factor 0.9--1.1)
  {[}CITATION{]}.
  \item \textbf{Spectrograms} receive a small circular time shift of up to
  $\pm 10$ frames and SpecAugment-style masking with two random time masks of up
  to 20 frames and two frequency masks of up to 6 bands, filled with the mean
  value of the spectrogram {[}CITATION{]}.
\end{itemize}

Table~\ref{tab:small-models} compares the standard and the small variants.

\begin{table}[htbp]
\centering
\small
\renewcommand{\arraystretch}{1.2}
\begin{tabularx}{\textwidth}{|p{2.6cm}|X|X|}
\hline
\textbf{Model} & \textbf{Standard} & \textbf{Small} \\
\hline
Early fusion (concatenation) & Conv1D 64 / 128; 49\,097 parameters & Conv1D 32 / 64; 18\,729 parameters \\
IMU backbone & Conv1D 64 / 128; 36\,297 parameters & Conv1D 32 / 64, L2, dropout 0.4; 12\,329 parameters \\
Audio backbone & Conv2D 32 / 64 with max pooling; 23\,945 parameters & Conv2D 16 / 32 with strides 4 / 2, L2, dropout 0.5; 7\,689 parameters \\
Label smoothing & -- & 0.1 (backbones) \\
Augmentation & -- & one augmented copy per training window (backbones) \\
\hline
\end{tabularx}
\caption{Standard and deployment-oriented (small) model variants. Parameter counts
refer to the nine-class output layer.}
\label{tab:small-models}
\end{table}

\section{Training Procedure}
\label{sec:training}

\subsection{Software and Hardware Environment}
\label{sec:environment}

All models were implemented in Python~3.10 with TensorFlow~2.10 and Keras
{[}CITATION{]}. Training ran on an NVIDIA GeForce
RTX~4070 Laptop GPU through the DirectML backend for Windows. This backend
supports TensorFlow only up to version~2.10, which fixed the framework version for
the whole project. Two parts of the pipeline could not run on the GPU with this
backend and had to be adapted: the short-time Fourier transform used for the
spectrograms was moved to the CPU, and the fused recurrent kernel of the GRU layer
was replaced by an unrolled implementation. Both changes affect only where and how
the computation is executed, not its result. Data preparation used NumPy, pandas
and librosa {[}CITATION{]}, and the classical baseline
used scikit-learn {[}CITATION{]}.

\subsection{Training Settings}
\label{sec:settings}

All neural networks were trained with the settings listed in
Table~\ref{tab:training}. The class imbalance described in
Section~\ref{sec:stats} is handled through class weights in the loss function.
The weight of each class is inversely proportional to its frequency in the
training split, so that errors on rare classes count more than errors on frequent
ones {[}CITATION{]}. The training data themselves are left unchanged, without
oversampling or undersampling.

\begin{table}[htbp]
\centering
\small
\renewcommand{\arraystretch}{1.2}
\begin{tabularx}{\textwidth}{|p{3.6cm}|X|}
\hline
\textbf{Setting} & \textbf{Value} \\
\hline
Optimiser & Adam {[}CITATION{]} \\
Learning rate & $10^{-3}$, constant \\
Batch size & 32 \\
Loss & Categorical cross-entropy (label smoothing 0.1 for the small backbones) \\
Class weighting & Inverse class frequency (``balanced'') \\
Epochs & 150 for backbones and early-fusion models; at most 80 for the fusion head \\
Model selection & Weights of the epoch with the highest validation accuracy \\
Fusion head & Early stopping (patience 15), learning rate halved on plateau (patience 6, minimum $10^{-5}$) \\
\hline
\end{tabularx}
\caption{Training settings shared by all neural network models.}
\label{tab:training}
\end{table}

The backbones and the early-fusion models are trained for a fixed 150 epochs at a
constant learning rate. On a small validation set, the validation accuracy can
fluctuate noticeably from one epoch to the next, so the weights at the end of
training are not necessarily the best ones. After training, each model is
therefore set back to the weights of the epoch with the highest validation
accuracy.

In the Keras version used here, the standard early-stopping callback restores the
best weights only when it actually stops training early. If training runs the full
number of epochs, the weights of the last epoch are kept without any warning. A
small custom callback was therefore written that stores the best weights during
training and restores them at the end. The effect of this choice is examined in
Chapter~\ref{ch:results}. Because the best epoch is selected on the validation
data, the procedure introduces a slight optimistic bias into the reported results,
a point taken up in the discussion of limitations in
Chapter~\ref{ch:discussion}.

\section{Evaluation Protocol}
\label{sec:evaluation}

\subsection{Grouped Train/Validation Split}
\label{sec:split}

The windows were created with 50\,\% overlap (Section~\ref{sec:windowing}), so
neighbouring windows share half of their samples. A random split would place such
neighbours on both sides of the train/validation boundary. The model would then
be evaluated on data that it had partly seen during training, and the measured
accuracy would overestimate its performance on new data {[}CITATION{]}.

To prevent this, the windows are first grouped into activity instances.
Consecutive windows with the same label belong to the same group as long as their
start times are at most 1.5\,s apart, that is, 1.5 times the window stride.
Consecutive windows of one uninterrupted activity occurrence start exactly 1\,s
apart, so the additional 0.5\,s serves as a margin: small irregularities in the
window start times do not split one occurrence into several groups, while
occurrences separated by a pause or by another activity, which leave a gap of at
least one further stride, are still kept apart. The 2565 windows form 125 such groups. The split is then made per class: the groups of
a class are shuffled, and whole groups are assigned to the validation set until it
holds at least 20\,\% of that class's windows. The remaining groups form the
training set. Every activity instance therefore lies entirely in one of the two
sets, and a fixed random seed makes the split identical for all models, so that
their results can be compared directly. Figure~\ref{fig:grouped-split} illustrates
the principle.

\begin{figure}[htbp]
\centering
\resizebox{\textwidth}{!}{%
\begin{tikzpicture}[x=1cm, y=1cm, font=\scriptsize]
  % activity instances on the timeline
  \foreach \a/\b/\s/\t in {0/3.5/10/{instance 1 (A)}, 3.5/6.5/30/{instance 2 (B)},
                           6.5/9.5/10/{instance 3 (A)}, 9.5/12/30/{instance 4 (B)}} {
    \fill[black!\s] (\a,0) rectangle (\b,0.6);
    \draw (\a,0) rectangle (\b,0.6);
    \node at ({(\a+\b)/2},0.3) {\t};
  }
  \draw[-{Latex[length=1.6mm]}] (0,-0.3) -- (12,-0.3) node[right] {time};
  % overlapping windows of instance 1
  \foreach \i in {0,1,2,3} \draw[|-|] (0.5*\i,-0.6-0.22*\i) -- (0.5*\i+2,-0.6-0.22*\i);
  \node[anchor=west] at (3.8,-0.93) {overlapping windows within instance 1};
  % grouped assignment
  \node[anchor=west, font=\scriptsize\bfseries] at (0,-1.75) {Grouped split (used):};
  \foreach \a/\b/\s in {0/3.5/train, 3.5/6.5/val, 6.5/9.5/val, 9.5/12/train}
    \draw[thick] (\a,-2.35) rectangle (\b,-2.0) node[midway] {\s};
  % random assignment
  \node[anchor=west, font=\scriptsize\bfseries] at (0,-2.85) {Random split (avoided):};
  \foreach \i/\s in {0/train,1/val,2/train,3/val,4/train}
    \draw (0.7*\i,-3.45) rectangle (0.7*\i+0.7,-3.1) node[midway] {\s};
  \node[anchor=west] at (3.8,-3.28) {overlapping windows of instance 1 fall on both sides};
\end{tikzpicture}}
\caption{Grouped splitting. Overlapping windows are grouped into activity
instances, and whole instances are assigned to either the training or the
validation set. A random split would distribute the overlapping windows of one
instance across both sets.}
\label{fig:grouped-split}
\end{figure}

Table~\ref{tab:split} shows the resulting split. Because only whole groups can be
moved, the validation share of a class is not exactly 20\,\%. It exceeds the
target considerably when a class consists of a few long instances, as is the case
for \texttt{walk}, \texttt{measuring} and \texttt{place\_container}, where a single
group already contains a large part of the class. Overall, 946 windows (37\,\%)
fall into the validation set and 1619 windows (63\,\%) into the training set. For
\texttt{place\_container}, only 23 windows remain for training. This uneven split
is a consequence of keeping activity instances intact, and its effect on the
results is discussed in Chapter~\ref{ch:discussion}.

\begin{table}[htbp]
\centering
\small
\renewcommand{\arraystretch}{1.15}
\begin{tabular}{|l|r|r|r|r|}
\hline
\textbf{Class} & \textbf{Training} & \textbf{Validation} & \textbf{Total} & \textbf{Validation share} \\
\hline
\texttt{assemble\_parts}   & 130 &  36 & 166 & 22\,\% \\
\texttt{carry\_container}  & 175 &  45 & 220 & 20\,\% \\
\texttt{hammering}         & 170 &  80 & 250 & 32\,\% \\
\texttt{idle}              & 187 & 120 & 307 & 39\,\% \\
\texttt{insert\_screws}    & 185 &  61 & 246 & 25\,\% \\
\texttt{measuring}         & 159 & 161 & 320 & 50\,\% \\
\texttt{place\_container}  &  23 &  51 &  74 & 69\,\% \\
\texttt{screw\_tightening} & 454 & 142 & 596 & 24\,\% \\
\texttt{walk}              & 136 & 250 & 386 & 65\,\% \\
\hline
\textbf{Total}             & \textbf{1619} & \textbf{946} & \textbf{2565} & \textbf{37\,\%} \\
\hline
\end{tabular}
\caption{Training and validation windows per class after the grouped split.}
\label{tab:split}
\end{table}

\subsection{Metrics}
\label{sec:metrics}

Overall accuracy, the share of correctly classified validation windows, is used as
the main measure for comparing the models. Because the classes are imbalanced,
accuracy alone can hide poor performance on small classes. Precision, recall and
F1-score are therefore also reported for each class, together with the
macro-averaged F1-score, which gives every class the same weight {[}CITATION{]}. Confusion matrices are used to identify which activities are mistaken
for one another. Unless stated otherwise, all results are reported on the
validation set of the grouped split and come from a single training run per
model.

\section{Chapter Summary}
\label{sec:ch4-summary}

This chapter introduced the models compared in this thesis. Inertial windows are
processed as $200 \times 6$ time series by a 1D CNN, and audio windows as
$198 \times 40$ log-mel spectrograms by a 2D CNN, for which three pooling variants
were developed to make use of the temporal structure of the sounds. The two
modalities are combined at the input level (channel concatenation and unified
spectrogram), at the level of learned embeddings, and at the level of class
probabilities, where a per-class weighting reflects the different reliability of
the two modalities across activities. A random forest on hand-crafted features
serves as a classical baseline, and reduced, regularised variants prepare the
deployment on the device. All models share the same input normalisation, class
weighting, training settings, best-epoch selection and grouped train/validation
split. Chapter~\ref{ch:implementation} describes how selected models were
converted for execution on the wearable device, and Chapter~\ref{ch:results}
presents the results.
